% An appendix
%======================================================================

\part*{APPENDICES}

\chapter{Proof of Concept – Prototyp Gry}
\label{appA}

Niniejszy dodatek zawiera opis funkcjonalnego prototypu gry \textit{Echoes of Vanitas} oraz dokumentację działających mechanik w formie opisowej. W prototypie zaimplementowane są następujące systemy:

\section{Działające Mechaniki Prototypu}

\subsection{Poruszanie się i Walka}
\begin{itemize}
    \item Gracz porusza się w czterech kierunkach klawiszami WASD,
    \item Atak wykonywany jest klawiszem Spacji,
    \item Dash (unik) dostępny jest klawiszem Shift,
    \item Trafienie gracza przez wroga zmniejsza punkty życia i może wywołać stan śmierci.
\end{itemize}

\subsection{System Przeciwników}
\begin{itemize}
    \item Przeciwnicy patrolują wyznaczone trasy (Path3D),
    \item Po wykryciu gracza przechodzą w stan pościgu,
    \item W zasięgu ataku wykonują obrażenia,
    \item Po trafieniu przez gracza wchodzą w stan ogłuszenia,
    \item Po śmierci odtwarzają animację, spawnują orby z przedmiotami i znikają z planszy,
    \item Licznik wrogów na ekranie HUD aktualizuje się w czasie rzeczywistym.
\end{itemize}

\subsection{System Questów}
\begin{itemize}
    \item NPC dający quest obecny jest w scenie; po wejściu w zasięg i wciśnięciu E quest jest przyjmowany,
    \item W menu questów (klawisz I → zakładka Questy) wyświetlają się aktywne zadania,
    \item Kliknięcie na quest wyświetla jego tytuł, opis, informacje o lokacji i nagrodach,
    \item Po zabiciiu wymaganej liczby wrogów quest jest automatycznie kończony,
    \item Po ukończeniu nagrody trafiają do inventory gracza.
\end{itemize}

\subsection{System Ekwipunku}
\begin{itemize}
    \item Przedmioty podniesione z podłogi trafiają do inventory (menu I → zakładka Ekwipunek),
    \item Gracz może przeciągać przedmioty między slotami inventory a slotami wyposażenia,
    \item Założenie przedmiotu zmienia statystyki postaci (widoczne w HUD),
    \item Przedmioty można łączyć w slotach Merge A+B i ulepszać przyciskiem UPGRADE,
    \item Slot Trash pozwala niszczyć niechciane przedmioty.
\end{itemize}

\section{Parametry Techniczne Prototypu}

\begin{itemize}
    \item \textbf{Silnik}: Godot 4.x z renderowaniem Forward Plus,
    \item \textbf{Język}: C\# (.NET 8),
    \item \textbf{Rozdzielczość}: 1920×1080,
    \item \textbf{Tryb renderowania}: 3D (środowisko 2.5D),
    \item \textbf{Platforma docelowa}: Windows (PC),
    \item \textbf{Liczba zaimplementowanych klas}: 59 plików \texttt{.cs},
    \item \textbf{Liczba scen Godot}: ponad 20 plików \texttt{.tscn}.
\end{itemize}

\section{Znane Ograniczenia Prototypu}

\begin{itemize}
    \item System dźwięku i muzyki nie jest jeszcze zaimplementowany,
    \item Pełny system rzutów kością opisany w założeniach fabularnych czeka na implementację,
    \item Dialogi z NPC ograniczają się do przyjmowania questów; system rozbudowanych rozmów (drzewo dialogowe) nie jest gotowy,
    \item Zapis stanu gry (save/load) nie jest zaimplementowany,
    \item Tylko jedna lokacja jest dostępna w prototypie.
\end{itemize}
